\documentclass{report}
\usepackage{amsmath,amssymb}
\setlength{\parindent}{0mm}
\setlength{\parskip}{1em}
\begin{document}
\begin{center}
\rule{6in}{1pt} \
{\large Amber, S. Robertson \\
{\bf Approximation of the Scattering Amplitude using Nonsymmetric Saddle Point Matrices}}

104 Mary Circle Hattiesburg \\ Ms 39402
\\
{\tt amber.sumner@eagles.usm.edu}\\
James, V. Lambers\end{center}

In this paper we seek to develop iterative methods that are more robust
and efficient than existing methods for solving the forward ($A{\bf
x}={\bf b}$) and adjoint ($A^{T}{\bf y}={\bf g}$) systems of linear
equations where the coefficient matrix is large, sparse, and
nonysmmetric. We can use these methods to approximate the scattering
amplitude defined by ${\bf g}^{T}{\bf x}$. We use a conjugate
gradient-like iteration for a nonsymmetric saddle point matrix that is
constructed so as to have a real positive spectrum, and a full set of
eigenvectors that are orthogonal in some sense. The preservation of these
properties allows us to generalize the conjugate gradient method to one
for nonsymmetric matrices. We find that this method performs more
consistently in the initial iterations than known methods for computing
the scattering amplitude such as GLSQR or QMR. We also look into using
techniques from "matrices, moments, and quadrature" to approximate
expressions of the form ${\bf u}^{T} f(A) {\bf v}$, where $A$ is a
nonsymmetric saddle point matrix, such as the scattering amplitude,
without solving the system directly. Block approaches and preconditioning
are also examined.


\end{document}
